%=======================================================================
%  4  SUSTAINABILITY & LIFE'S PRINCIPLES
%=======================================================================
\section{Sustainability \& Life's Principles}
\label{sec:sustainability}

\subsection{Life's Principles applied to Concept V2}
\label{sec:lp_applied}

\subsubsection{Be resource efficient --- minimize the use of external resources}
\label{sec:lp_resource}

Cooling needs no auxiliary electricity: circulation is driven by buoyancy alone
(Section~\ref{sec:check1}), and the loop is closed, so no water is consumed in operation
(Table~\ref{tab:kpicomp}). Pumpless operation is not demonstrated. The mains alone use
most of the $\approx$300\,Pa buoyancy pressure, the drawn roof coils exceed it by a factor
of about 67, and refilling after drainback may need a pump (Section~\ref{sec:v2check}).

\subsubsection{Adapt to changing conditions --- day/night and seasonal cycles}
\label{sec:lp_adapt}

The daily cycle is used rather than fought: a hotter façade produces a larger density
difference and more flow, so cooling rises and falls with the load without a controller
(Section~\ref{sec:check1}), and the roof releases most of its heat to the night sky
(Section~\ref{sec:check2}). The seasonal cycle is met by draining the loop before frost
and refilling it in spring (Section~\ref{sec:v2check}). The seasonal adaptation therefore
functions as a withdrawal rather than an adaptation. Even after draining the pipes there
is a remaining risk of freezing leftover water. Furthermore the ribs and the reflective
roof act the same even in winter.

\subsubsection{Use life-friendly chemistry --- no toxic chemicals or materials}
\label{sec:lp_chemistry}

The heat-transfer fluid is plain water. Frost protection is solved by draining the pipes
instead of antifreeze such as glycol. Reflection and emission are sought through the
surface optics of the silver ant \cite{shi2015}, not through a functional additive. We
meet this principle by leaving substances out; we have not verified the ones we use. Two
points remain open and unchecked: the timber treatment of the ribs and the binder and
preservatives of the roof coating.

\subsection{Environmental benefits, limitations and trade-offs}
\label{sec:benefits}

\subsubsection{Environmental benefits}

The climate benefit is what V2 replaces: every hour of air conditioning it makes
unnecessary is avoided emissions and avoided waste heat in the street
(Section~\ref{sec:problem}). Radiation to the night sky carries the collected heat out
of the urban air. The other release options of Table~\ref{tab:morph}, exchange with the
air beside the wall or evaporation, leave it in the street (Section~\ref{sec:check2}).
The street itself should be cooler. The shaded façade radiates less onto the pavement,
there is no condenser at pedestrian height, and reflection happens on the roof, where it
does not raise heat stress in the canyon the way a reflective wall does \cite{lee2018}.
The retrofit is reversible, nothing is demolished. The panels can be removed and the
building returned to its previous state, so the retrofit creates no demolition waste and
consumes no land. In operation the system consumes no water and contains no refrigerant;
what it does release is the abrasion of the roof coating into the run-off
(Section~\ref{sec:materials}).

\subsubsection{Limitations}

The export of heat depends on a clear, dry night sky. In the humid, hazy nights of a heat
wave most of the heat only warms the roof air, so the concept helps least in the
conditions it is meant for (Section~\ref{sec:check2}). Outside the summer the drained
loop does nothing, and no other environmental function fills the gap. No cooling figure
exists, so the avoided emissions cannot be counted, and the whole benefit rests on
circulation without a pump, the first open check of Table~\ref{tab:priorities}. The added
materials carry embodied emissions that the constraint in Table~\ref{tab:constraints}
asks to be repaid within the service life; we carried out no life cycle assessment, so
neither side of that comparison exists.

\subsubsection{Trade-offs}

Ribs and panels change how the building looks. On a listed façade this collides with
heritage protection \cite{stadtzuerich_web}, \cite[pp.~148--151]{stadtzuerich2020}. What
shades and reflects in summer also rejects winter sun and can raise heating demand and
its emissions \cite{epa_coolroofs}. The net annual effect is unknown. Every square meter
of roof that releases heat is a square meter without photovoltaic yield
(Figure~\ref{fig:roof}). Avoiding antifreeze costs the shoulder seasons and admits air
that may bring back a corrosion inhibitor (Section~\ref{sec:v2check}). Repairable
external pipes are the same visible layer that alters the façade.

\subsection{Materials, durability, repair and circularity}
\label{sec:materials}

The timber ribs are worn by UV, driving rain and freeze--thaw, so they belong to the
biological cycle \cite{braungart2002}. What weathers off lands on the pavement, and the
timber must be untreated, naturally durable, or protected by detailing rather than by
chemistry. The collector panels and their finish see slow surface wear and are otherwise
used as structure: technical cycle, if the panel is one recoverable material with a
benign finish. The panels are designed to be removable (Figure~\ref{fig:v2pipes}). The
roof coating is worn by UV, rain, soiling and cleaning, and abraded into run-off that
reaches the stormwater network of Section~\ref{sec:systemview}, so it must break down
harmlessly. Pipework, carrier plate, supports and fixings are used, not consumed:
technical cycle, recoverable if specified as separable single materials. The water in the
loop is used but mobile and fits neither cycle cleanly. About 5.8\,m\textsuperscript{3} is
drained to storage each winter and eventually disposed of, and with an inhibitor it is
the one component that can leave the system by accident. Measured against the eight
cycle design principles of the lecture \cite{zhaw2026_c2c}, V2 answers four in principle
and none with evidence. Modular construction (3) and reparability (5) follow from moving
the pipes out of the wall: a panel can be removed and the façade pipes inspected without
opening the wall (Sections~\ref{sec:review} to \ref{sec:v2access}). Ecotoxic chemicals
(1) and recyclable materials (2) are answered above.

\runin{Material checked so far.} White cool-roof coatings owe their reflectance to
titanium dioxide, so it was checked in the REACH database of the European Chemicals
Agency \cite{echa_reach}. The powder form was classified in 2020 as a suspected
carcinogen by inhalation (Regulation (EU) 2020/217, H351). The General Court annulled the
classification in 2022 and the Court of Justice confirmed the annulment on 1~August 2025
\cite{cjeu2025}, so the pigment currently carries no harmonized hazard classification.
The result is not clear. The concern was dust inhalation, not the applied film, and the
ruling turned on the reliability of one study, not on a finding of safety. The pigment is
also not the used coating. The binder and the film preservatives are what reach the
run-off instead. One substance of at least four has been checked. The other three belong
to item~5 of Table~\ref{tab:priorities}.

\clearpage

\subsection{Did sustainability considerations change the design?}
\label{sec:designchanges}

The design was adjusted in four instances which were documented in
Sections~\ref{sec:creating} and \ref{sec:evaluation}.

\runin{1. The reflective surface moved from the façade to the roof.} Concept~C reflected
on the façade and was set aside because reflective walls raised pedestrian heat stress in
simulations of a Stuttgart street canyon \cite{lee2018}; Concept~A reflects on the roof
instead (Section~\ref{sec:morph}). The change was made for the people at street level,
not for cooling performance.

\runin{2. Water and power demand removed two options before any concept was drawn.} The
evaporating wet layer and the fan-driven cavity of Table~\ref{tab:morph} were not
developed further because of their water and power demand (Section~\ref{sec:morph}). The
Energy and Water constraints of Table~\ref{tab:constraints}, which are the principle
``be resource efficient'' written as constraints, decided this.

\runin{3. The pipes moved out of the wall.} V1 embedded them in the façade. The
engineering review found that in case of damage the lack of access was unacceptable
(Table~\ref{tab:review}). V2 therefore placed them in separable external panels with a
drainage gap behind them (Section~\ref{sec:v2}).

\runin{4. Water instead of glycol.} Frost protection was decided by drainback rather
than by an additive (Section~\ref{sec:v2check}).

\subsection{Limits of this assessment}
\label{sec:limits}

The assessment is made against Concept V2 as described in Section~\ref{sec:v2}, whose
central claim of buoyancy-only circulation is the first-ranked open check of
Table~\ref{tab:priorities}. If a pump is needed, the Energy constraint and the
resource-efficiency assessment of Section~\ref{sec:lp_applied} fall with it. Cascade use
and return logistics lie outside the project boundary, and of the substances the concept
would put into the street one has been checked and three have not.

\clearpage
